% =====================================================================
%  2.3 Andregradslikninger — teori og oppgaver
% =====================================================================

\del{1}{Faktorisering og nullproduktregelen}

\begin{ordliste}
andregradslikning & En likning der den høyeste potensen av $x$ er $x^2$. Eksempel: $x^2 - 4x + 3 = 0$. & \mittsprak\\
nullproduktregelen & Hvis $a \cdot b = 0$, er $a = 0$ eller $b = 0$. & \mittsprak\\
sum--produkt & Metode: finn to tall med gitt produkt og gitt sum, og faktoriser $x^2 + bx + c$. & \mittsprak\\
løsning & En verdi av $x$ som gjør likningen sann. En andregradslikning har 0, 1 eller 2 løsninger. & \mittsprak\\
\end{ordliste}

\begin{mustkunne}
\textbf{Nullproduktregelen.} Hvis $a \cdot b = 0$, er $a = 0$ eller $b = 0$.
\emph{Hvorfor?} Hvis $a \neq 0$, kan vi dele begge sider på $a$. Da blir $b = 0$.
Regelen gjelder \textbf{bare for 0}. Fra $a \cdot b = 6$ kan vi ikke si noe om $a$ eller $b$.

\smallskip
\textbf{Se på likningen først.} Flytt alt over på én side, slik at høyre side er 0. Velg så metode:

\smallskip
\renewcommand{\arraystretch}{1.5}
\begin{tabular}{@{}l l l@{}}
\textbf{Felles faktor} & $x^2 - 5x = 0 \;\Rightarrow\; x(x-5) = 0$ & $x = 0$ eller $x = 5$\\
\textbf{Kvadratrot} & $x^2 = 9$ & $x = 3$ eller $x = -3$\\
\textbf{Kvadratsetning} & $x^2 - 6x + 9 = 0 \;\Rightarrow\; (x-3)^2 = 0$ & $x = 3$\\
\textbf{Sum--produkt} & $x^2 + 5x + 6 = (x+2)(x+3)$ & $x = -2$ eller $x = -3$\\
\end{tabular}

\smallskip
Sum--produkt: finn to tall med produkt 6 og sum 5. Metoden passer best når $a = 1$ og løsningene er hele tall.
\end{mustkunne}

\begin{eksempel}[1 — felles faktor, ikke del på $x$]
Løs likningen $3x^2 = 12x$.

\begin{stegtabell}
$3x^2 - 12x = 0$ & Flytt $12x$ over, slik at høyre side er 0.\\
$3x(x - 4) = 0$ & Sett $3x$ utenfor parentes. Vi deler ikke på $x$.\\
$3x = 0$ eller $x - 4 = 0$ & Nullproduktregelen.\\
$x = 0$ eller $x = 4$ & Løs hver av dem.\\
Kontroll: & $x = 4$ gir $3 \cdot 16 = 48$ og $12 \cdot 4 = 48$. $x = 0$ gir $0 = 0$.\\
\end{stegtabell}
\end{eksempel}

\begin{eksempel}[2 — sum--produkt]
Løs likningen $x^2 - x - 12 = 0$.

\begin{stegtabell}
Produkt $-12$, sum $-1$ & Tallene er $-4$ og $3$, fordi $(-4) \cdot 3 = -12$ og $-4 + 3 = -1$.\\
$(x - 4)(x + 3) = 0$ & Faktoriser.\\
$x = 4$ eller $x = -3$ & Nullproduktregelen.\\
Kontroll: & $4^2 - 4 - 12 = 0$ og $(-3)^2 + 3 - 12 = 0$.\\
\end{stegtabell}
\end{eksempel}

\begin{vanligfeil}
\begin{itemize}[leftmargin=1.2em,itemsep=1pt,topsep=0pt]
  \item $x^2 = 5x \;\Rightarrow\; x = 5$ er feil. Når du deler på $x$, mister du $x = 0$. Riktig: $x(x - 5) = 0$, så $x = 0$ eller $x = 5$.
  \item $x^2 = 9 \;\Rightarrow\; x = 3$ er feil. Også $(-3)^2 = 9$. Riktig: $x = 3$ eller $x = -3$.
  \item $(x+1)(x+2) = 12 \;\Rightarrow\; x + 1 = 12$ eller $x + 2 = 12$ er feil. Regelen gjelder bare for 0. Riktig: gang ut til $x^2 + 3x - 10 = 0$, så $(x+5)(x-2) = 0$ og $x = -5$ eller $x = 2$.
\end{itemize}
\end{vanligfeil}

\nivaahode{1}

\oppg{1.1}{1}{\uh}
Løs likningene.
\begin{tasks}(2)
  \task $(x - 1)(x + 4) = 0$
  \task $x(x - 7) = 0$
  \task $(2x - 6)(x + 5) = 0$
  \task $3(x + 2)(x - 8) = 0$
\end{tasks}

\oppg{1.2}{1}{\uh}
Løs likningene. Skriv begge løsningene.
\begin{tasks}(4)
  \task $x^2 - 8x = 0$
  \task $2x^2 + 10x = 0$
  \task $x^2 = 49$
  \task $2x^2 = 32$
\end{tasks}

\oppg{1.3}{1}{\uh}
Vi skal løse $x^2 + 7x + 10 = 0$ med sum--produkt-metoden. Fyll inn.
\begin{tasks}(1)
  \task Tallene har produkt 10 og sum 7. De er \skrivlinje[1.5cm] og \skrivlinje[1.5cm].
  \task $x^2 + 7x + 10 = (x + $\skrivlinje[1.2cm]$)(x + $\skrivlinje[1.2cm]$)$
  \task $x = $\skrivlinje[1.5cm] eller $x = $\skrivlinje[1.5cm]
\end{tasks}

\nivaahode{2}

\oppg{1.4}{2}{\uh}
Løs likningene. Skriv hvilken metode du valgte, og hvorfor.
\begin{tasks}(2)
  \task $x^2 - 3x - 10 = 0$
  \task $x^2 - 10x + 25 = 0$
  \task $6x^2 = 18x$
  \task $x^2 + 2x = 15$
\end{tasks}

\oppg{1.5}{2}{\uh}
Løs likningene. Husk at nullproduktregelen bare gjelder når den ene siden er 0.
\begin{tasks}(2)
  \task $(x - 3)(x + 2) = 14$
  \task $x(x + 3) = 28$
\end{tasks}

\oppg{1.6}{2}{\uh}
Ola løser en likning slik:
\begin{align*}
&x^2 + 5x = 4x + 20\\
&\longrightarrow\quad x^2 + 5x - 4x - 20 = 0\\
&\longrightarrow\quad x^2 + x - 20 = 0\\
&\longrightarrow\quad (x + 5)(x - 4) = 0\\
&\longrightarrow\quad x = 5 \text{ eller } x = -4
\end{align*}
Finn den første feilen. Løs likningen riktig.

\nivaahode{3}

\oppg{1.7}{3}{\uh}
Er påstanden alltid, noen ganger eller aldri sann? Begrunn svaret.
\begin{tasks}(1)
  \task Likningen $x^2 = k$ har to løsninger.
  \task Likningen $(x - m)(x - n) = 0$ har to ulike løsninger.
  \task Tallene $0$ og $-b$ er løsninger av $x^2 + bx = 0$.
\end{tasks}

\oppg{1.8}{3}{\uh}
Finn to verdier av $k$ slik at $x^2 + kx + 36 = 0$ har nøyaktig én løsning.
\hint{Hvordan ser venstre side ut faktorisert når det bare er én løsning?}

\oppg{1.9}{3}{\uh}
En andregradslikning $x^2 + px + q = 0$ har løsningene $x = 3$ og $x = 8$.
\begin{tasks}(1)
  \task Finn $p$ og $q$.
  \task Hvordan henger $p$ og $q$ sammen med løsningene? Forklar med faktorisering.
\end{tasks}

\begin{sjekk}
\begin{tasks}(2)
  \task Løs $x^2 + 7x = 0$.
  \task Hvor mange løsninger har $x^2 = -16$?
  \task! Ola løser $(x-1)(x+2) = 4$ ved å sette $x - 1 = 4$ eller $x + 2 = 4$. Forklar hva som er feil.
\end{tasks}
\end{sjekk}

% ---------------------------------------------------------------------

\del{2}{abc-formelen og diskriminanten}

\begin{ordliste}
abc-formelen & En formel som løser alle andregradslikninger $ax^2 + bx + c = 0$. & \mittsprak\\
diskriminant & Tallet $D = b^2 - 4ac$. Det sier hvor mange løsninger likningen har. & \mittsprak\\
fullstendig kvadrat & Et uttrykk som kan skrives som $(x + d)^2$, for eksempel $x^2 + 6x + 9$. & \mittsprak\\
nullpunkt & En $x$-verdi der funksjonsverdien er 0. Der treffer grafen $x$-aksen. & \mittsprak\\
\end{ordliste}

\begin{mustkunne}
\textbf{Utledning med fullstendig kvadrat.} Vi lager et kvadrat på venstre side.

\smallskip
\begin{minipage}[t]{0.46\linewidth}
\textbf{Tall:} $x^2 + 6x - 7 = 0$
\begin{align*}
x^2 + 6x &= 7\\
x^2 + 6x + 3^2 &= 7 + 3^2\\
(x + 3)^2 &= 16\\
x + 3 &= \pm 4\\
x &= 1 \text{ eller } x = -7
\end{align*}
\end{minipage}\hfill
\begin{minipage}[t]{0.50\linewidth}
\textbf{Generelt:} $ax^2 + bx + c = 0$
\begin{align*}
x^2 + \frac{b}{a}x + \frac{c}{a} &= 0 \quad \text{(del på } a \neq 0)\\
x^2 + \frac{b}{a}x &= -\frac{c}{a}\\
x^2 + \frac{b}{a}x + \left(\frac{b}{2a}\right)^2 &= \left(\frac{b}{2a}\right)^2 - \frac{c}{a}\\
\left(x + \frac{b}{2a}\right)^2 &= \frac{b^2 - 4ac}{4a^2}\\
x + \frac{b}{2a} &= \pm \frac{\sqrt{b^2 - 4ac}}{2a}
\end{align*}
\end{minipage}

\medskip
\[ \boxed{\;x = \frac{-b \pm \sqrt{b^2 - 4ac}}{2a}\;} \qquad D = b^2 - 4ac \]
$D > 0$: to løsninger. $D = 0$: én løsning. $D < 0$: ingen løsning, fordi $\left(x + \frac{b}{2a}\right)^2$ ikke kan være negativ.

\medskip
\begin{minipage}[t]{0.32\linewidth}
\centering
\begin{tikzpicture}
\begin{axis}[width=\linewidth, height=0.8\linewidth, axis lines=middle, xmin=-2, xmax=4, ymin=-4, ymax=4,
  x tick label style={anchor=north east, font=\scriptsize}, xtick={-1,3}, ytick={\empty}, tick label style={font=\scriptsize}, xlabel={$x$}, ylabel={$y$},
  label style={font=\small}, title={\small $D > 0$: to nullpunkt}]
\addplot[domain=-1.5:3.5, thick, hoved, samples=40] {x^2 - 2*x - 3};
\addplot[only marks, mark=*, mark size=1.5pt] coordinates {(-1,0) (3,0)};
\end{axis}
\end{tikzpicture}
\end{minipage}\hfill
\begin{minipage}[t]{0.32\linewidth}
\centering
\begin{tikzpicture}
\begin{axis}[width=\linewidth, height=0.8\linewidth, axis lines=middle, xmin=-2, xmax=4, ymin=-4, ymax=4,
  xtick={1}, ytick={\empty}, tick label style={font=\scriptsize}, xlabel={$x$}, ylabel={$y$},
  label style={font=\small}, title={\small $D = 0$: ett nullpunkt}]
\addplot[domain=-1:3, thick, hoved, samples=40] {x^2 - 2*x + 1};
\addplot[only marks, mark=*, mark size=1.5pt] coordinates {(1,0)};
\end{axis}
\end{tikzpicture}
\end{minipage}\hfill
\begin{minipage}[t]{0.32\linewidth}
\centering
\begin{tikzpicture}
\begin{axis}[width=\linewidth, height=0.8\linewidth, axis lines=middle, xmin=-2, xmax=4, ymin=-4, ymax=4,
  xtick={1}, ytick={\empty}, tick label style={font=\scriptsize}, xlabel={$x$}, ylabel={$y$},
  label style={font=\small}, title={\small $D < 0$: ingen nullpunkt}]
\addplot[domain=-0.5:2.5, thick, hoved, samples=40] {x^2 - 2*x + 3};
\end{axis}
\end{tikzpicture}
\end{minipage}

\medskip
\textbf{Faktorisering.} Har $ax^2 + bx + c = 0$ løsningene $x_1$ og $x_2$, er
$ax^2 + bx + c = a(x - x_1)(x - x_2)$. Løsningen $x_1 = 2$ gir faktoren $(x - 2)$, ikke $(x + 2)$.
\end{mustkunne}

\begin{eksempel}[1 — abc-formelen med kontroll]
Løs likningen $2x^2 - 3x - 2 = 0$.

\begin{stegtabell}
$a = 2,\; b = -3,\; c = -2$ & Les av tallene. Ta med fortegnet.\\
$x = \dfrac{-(-3) \pm \sqrt{(-3)^2 - 4 \cdot 2 \cdot (-2)}}{2 \cdot 2}$ & Sett inn i abc-formelen. Bruk parentes rundt $b$.\\
$x = \dfrac{3 \pm \sqrt{9 + 16}}{4} = \dfrac{3 \pm 5}{4}$ & $-(-3) = 3$, $(-3)^2 = 9$ og $-4 \cdot 2 \cdot (-2) = 16$. $D = 25$ og $\sqrt{25} = 5$.\\
$x = 2$ eller $x = -\dfrac{1}{2}$ & $\frac{3+5}{4} = 2$ og $\frac{3-5}{4} = -\frac{1}{2}$.\\
Kontroll: & $2 \cdot 2^2 - 3 \cdot 2 - 2 = 0$ og $2 \cdot \frac14 + \frac32 - 2 = 0$.\\
\end{stegtabell}
\end{eksempel}

\begin{eksempel}[2 — antall løsninger uten å løse]
Avgjør hvor mange løsninger hver likning har.

\begin{stegtabell}
$x^2 - 4x + 7 = 0$ & $D = (-4)^2 - 4 \cdot 1 \cdot 7 = 16 - 28 = -12$. $D < 0$: ingen løsning.\\
$2x^2 + x - 3 = 0$ & $D = 1^2 - 4 \cdot 2 \cdot (-3) = 1 + 24 = 25$. $D > 0$: to løsninger.\\
$4x^2 + 12x + 9 = 0$ & $D = 12^2 - 4 \cdot 4 \cdot 9 = 144 - 144 = 0$. $D = 0$: én løsning.\\
\end{stegtabell}
\end{eksempel}

\begin{vanligfeil}
\begin{itemize}[leftmargin=1.2em,itemsep=1pt,topsep=0pt]
  \item $b = -3$ gir $-b^2 = -9$. Det er feil. Det er $b^2 = (-3)^2 = 9$ under rottegnet, og $-b = 3$ foran.
  \item $x = -b \pm \dfrac{\sqrt{D}}{2a}$ er feil. Hele telleren deles på $2a$: $x = \dfrac{-b \pm \sqrt{D}}{2a}$.
  \item I $x^2 + 2x = 8$ er $c$ ikke 8. Flytt først: $x^2 + 2x - 8 = 0$ gir $c = -8$.
  \item Nullpunktene $2$ og $-3$ gir $(x - 2)(x + 3)$, ikke $(x + 2)(x - 3)$.
\end{itemize}
\end{vanligfeil}

\nivaahode{1}

\oppg{2.1}{1}{\uh}
Regn ut diskriminanten $D = b^2 - 4ac$. Si hvor mange løsninger likningen har.
\begin{tasks}(3)
  \task $x^2 + 2x + 5 = 0$
  \task $3x^2 - 2x - 1 = 0$
  \task $4x^2 - 4x + 1 = 0$
\end{tasks}

\oppg{2.2}{1}{\uh}
Løs likningene med abc-formelen. Gi eksakte svar.
\begin{tasks}(3)
  \task $x^2 - 7x + 10 = 0$
  \task $x^2 + 2x - 4 = 0$
  \task $2x^2 + 5x + 2 = 0$
\end{tasks}

\oppg{2.3}{1}{\uh}
Løs $x^2 + 4x - 21 = 0$ med fullstendig kvadrat. Fyll inn.
\begin{tasks}(1)
  \task $x^2 + 4x = $\skrivlinje[1.2cm]. Legg til $2^2 = 4$ på begge sider: $x^2 + 4x + 4 = $\skrivlinje[1.2cm]
  \task $(x + 2)^2 = $\skrivlinje[1.2cm], så $x + 2 = \pm$\skrivlinje[1.2cm], altså $x = $\skrivlinje[1.2cm] eller $x = $\skrivlinje[1.2cm]
\end{tasks}

\nivaahode{2}

\oppg{2.4}{2}{\uh}
Løs likningene. Skriv hvilken metode du valgte, og hvorfor. Gi eksakte svar.
\begin{tasks}(2)
  \task $x^2 - 2x - 2 = 0$
  \task $3x^2 + x - 2 = 0$
  \task $2x^2 - 6x + 5 = 0$
  \task $5x^2 - 20 = 0$
\end{tasks}

\oppg{2.5}{2}{\uh}
Emma løser en likning slik:
\begin{align*}
&2x^2 - 5x - 3 = 0, \quad a = 2,\; b = -5,\; c = -3\\
&\longrightarrow\quad x = \frac{5 \pm \sqrt{(-5)^2 - 4 \cdot 2 \cdot (-3)}}{2 \cdot 2}\\
&\longrightarrow\quad x = \frac{5 \pm \sqrt{-25 + 24}}{4}\\
&\longrightarrow\quad \text{Ingen løsning, fordi } D = -1 < 0
\end{align*}
Finn den første feilen. Løs likningen riktig.

\oppg{2.6}{2}{\mh}
Et rektangel har omkrets 34 cm og areal 60 cm$^2$. Kall den ene siden $x$. Sett opp en likning for arealet og finn sidelengdene.

\oppg{2.7}{2}{\mh}
En ball kastes rett opp. Høyden over bakken er $h(t) = 20t - 5t^2$ meter, der $t$ er tiden i sekunder. Finn når ballen er 15 meter over bakken. Forklar hva de to svarene betyr.

\nivaahode{3}

\oppg{2.8}{3}{\uh}
Se på parabelen $y = x^2 - 4x + k$. Finn \textbf{alle} verdier av $k$ slik at parabelen
\begin{tasks}(3)
  \task skjærer $x$-aksen i to punkt
  \task tangerer $x$-aksen
  \task ikke treffer $x$-aksen
  \task! Bruk diskriminanten til å vise at $k = \pm 12$ er de eneste verdiene i oppgave 1.8.
\end{tasks}

\oppg{2.9}{3}{\uh}
Gitt likningen $ax^2 + bx + c = 0$. Er påstanden alltid, noen ganger eller aldri sann? Begrunn svaret.
\begin{tasks}(1)
  \task Har $a$ og $c$ motsatt fortegn, har likningen to løsninger.
  \task Har $a$ og $c$ samme fortegn, har likningen ingen løsning.
\end{tasks}

\oppg{2.10}{3}{\uh}
Grafen til $y = ax^2 + bx + c$ skjærer $x$-aksen i $-1$ og $3$, og $y$-aksen i $6$. Finn $a$, $b$ og $c$.

\begin{sjekk}
\begin{tasks}(2)
  \task Løs $2x^2 - x - 1 = 0$.
  \task Hvor mange løsninger har $x^2 + x + 1 = 0$?
  \task! Hva forteller $D = 0$ om grafen til $y = ax^2 + bx + c$?
\end{tasks}
\end{sjekk}
